\section{Reproducibility Verification Market}

\label{sec:market}

\subsection{Market Design}

The verification market creates economic incentives for independent reproducibility testing.
The market operates through a three-phase mechanism:

\textbf{Phase 1: Bounty Posting.}
When an experiment reaches the \texttt{completed} status, a verification bounty is automatically created.
The bounty amount $B$ is determined by:

\begin{equation}
    B = B_{\text{base}} + B_{\text{complexity}} \cdot \log(|V|) + B_{\text{impact}} \cdot I(E)
\end{equation}

\noindent where $B_{\text{base}}$ is a minimum bounty funded by the experiment's staking deposit, $|V|$ is the number of computational units in the execution graph, and $I(E)$ is an impact score based on citations, downloads, and dependent experiments.

\textbf{Phase 2: Verification Execution.}
Verifiers claim bounties by re-executing the experiment's execution graph on independent infrastructure.
The verification is structured as:

\begin{equation}
    V_{\text{result}} = \begin{cases}
        \text{REPRODUCED} & \text{if } d(\text{output}_{\text{original}}, \text{output}_{\text{verified}}) \leq \epsilon \\
        \text{PARTIAL} & \text{if } \frac{|\{v_i : d_i \leq \epsilon_i\}|}{|V|} \geq 0.8 \\
        \text{FAILED} & \text{otherwise}
    \end{cases}
\end{equation}

\textbf{Phase 3: Dispute Resolution.}
When verification fails, a dispute resolution process convenes a panel of domain-qualified verifiers who independently assess whether the failure stems from the original experiment, the verification attempt, or environmental non-determinism.

\subsection{Incentive Analysis}

We analyze the incentive structure using mechanism design theory.
Let $c_v$ be the cost of verification and $B$ be the bounty.
A rational verifier participates when:

\begin{equation}
    \mathbb{E}[B \cdot P(\text{accept})] > c_v + c_{\text{opportunity}}
\end{equation}

The protocol ensures incentive compatibility through three mechanisms:

\begin{enumerate}[leftmargin=*]
    \item \textbf{Staking:} Both experimenters and verifiers stake tokens that are slashed for dishonest behavior.
    \item \textbf{Reputation:} Successful verifications increase a verifier's reputation score, granting access to higher-value bounties.
    \item \textbf{Insurance:} A protocol-level insurance pool compensates verifiers whose honest work is challenged through frivolous disputes.
\end{enumerate}

\subsection{Semantic Reproducibility Proofs}

We introduce Semantic Reproducibility Proofs (SRP) that go beyond bitwise output comparison to assess whether results are scientifically equivalent:

\begin{equation}
    \text{SRP}(E_1, E_2) = \bigwedge_{j=1}^{m} \text{Consistent}(\text{conclusion}_j(E_1), \text{conclusion}_j(E_2))
\end{equation}

\noindent where $\text{conclusion}_j$ extracts the $j$-th scientific conclusion (e.g., hypothesis test outcome, effect size direction, confidence interval overlap) and $\text{Consistent}$ evaluates agreement accounting for expected stochastic variation.

The SRP framework handles three common sources of acceptable variation:

\begin{table}[H]
\centering
\caption{Sources of acceptable variation in reproducibility assessment.}
\label{tab:variation}
\begin{tabularx}{\textwidth}{@{}llX@{}}
\toprule
\textbf{Source} & \textbf{Detection} & \textbf{Resolution} \\
\midrule
Floating-point non-determinism & Bitwise output diff & Statistical equivalence test on aggregated metrics \\
Stochastic algorithms & Seed-dependent paths & Multi-seed execution with distribution comparison \\
Platform differences & Hardware attestation diff & Cross-platform tolerance bounds from calibration runs \\
\bottomrule
\end{tabularx}
\end{table}

% ============================================================
% 7. Implementation
% ============================================================
